\documentclass{article}

% ICML 2025 style
\usepackage[accepted]{icml2025}

% Standard packages
\usepackage[utf8]{inputenc}
\usepackage[T1]{fontenc}
\usepackage{hyperref}
\usepackage{url}
\usepackage{booktabs}
\usepackage{amsmath}
\usepackage{amssymb}
\usepackage{graphicx}
\usepackage{microtype}
\usepackage{xcolor}
\usepackage{multirow}
\usepackage{enumitem}

\graphicspath{{./}}

\tolerance=800
\emergencystretch=12pt

% -----------------------------------------------------------------------
\icmltitlerunning{Cross-Lingual Hate Speech Detection: English to Urdu Transfer}

\begin{document}

\twocolumn[
\icmltitle{Cross-Lingual Hate Speech Detection:\\
           Fine-Tuning on English, Transferring to Urdu}

\begin{icmlauthorlist}
\icmlauthor{Muteeullah Baig}{uu}
\icmlauthor{Rana Muhammad Hamza}{uu}
\end{icmlauthorlist}

\icmlaffiliation{uu}{AI-600 Deep Learning, LUMS, Spring 2026}
\icmlcorrespondingauthor{Muteeullah Baig}{25280036@lums.edu.pk}
\icmlcorrespondingauthor{Rana Muhammad Hamza}{25280103@lums.edu.pk}

\icmlkeywords{hate speech detection, cross-lingual transfer, multilingual transformers,
              Urdu, XLM-RoBERTa, few-shot learning, ISE-Hate, fairness}

\vskip 0.3in
]

\printAffiliationsAndNotice{}
{\small \textbf{Code:}
\url{https://github.com/MuteeullahBaig/AI600-HateSpeech-CrossLingual}}


% -----------------------------------------------------------------------
\begin{abstract}
We study cross-lingual transfer of hate speech detection from English to
Urdu (Nastaliq script), using the ISE-Hate dataset of 21{,}759 annotated
Pakistani tweets~\citep{akram2023ise}.
Three multilingual transformers (mBERT, XLM-R-base, XLM-R-large) are
fine-tuned on the English Davidson et al.\ (2017) corpus, then evaluated
zero-shot and few-shot on Urdu.
Despite Urdu being one of XLM-R's 100 pre-training languages, zero-shot
hate~F1 reaches only 0.147, driven by class-imbalance bias absorbed from
English training data.
Unlike transfer to out-of-vocabulary scripts, few-shot adaptation does
\emph{not} collapse: hate~F1 rises steadily from 0.590 ($K=8$) to
0.694 ($K=64$), confirming that correct target-language selection is
critical for stable few-shot transfer.
Sub-category fairness analysis using ISE-Hate Level-2 labels reveals
a 20$\times$ disparity in zero-shot recall between Interfaith hate (28.2\%)
and generic ``Other'' hate (1.4\%), exposing systematic blind spots in the
English-calibrated model.
Classical TF-IDF~$+$~Logistic Regression on full Urdu supervision achieves
hate~F1~$=$~0.747, while supervised XLM-R-large fine-tuning provides the
neural upper bound.
We further discuss how annotator biases in the English training source
are directly transmitted into cross-lingual transfer behaviour.
\end{abstract}

% -----------------------------------------------------------------------
\section{Introduction}
\label{sec:intro}

Hate speech on social media disproportionately affects non-English-speaking
communities, yet most automated detection systems are built for English.
Urdu, spoken by over 230 million people and the national language of Pakistan,
is a prime example: it is widely used on social media but remains
substantially under-resourced for NLP.

This work investigates \textit{cross-lingual transfer learning} for hate
speech detection, treating English as the source language and Urdu as the
target, using the ISE-Hate corpus of Pakistani tweets written in Nastaliq
(Arabic) script~\citep{akram2023ise}.
Nastaliq Urdu is one of the 100 languages in XLM-RoBERTa's pre-training
data~\citep{conneau2020unsupervised}, making it a methodologically sound
target for evaluating cross-lingual transfer.

We address four research questions:

\begin{itemize}[leftmargin=*,nosep]
  \item \textbf{RQ1 (Backbone):} How does model capacity affect English
        detection and cross-lingual transfer to Urdu?
  \item \textbf{RQ2 (Zero-shot):} Can a model trained only on English detect
        Urdu hate speech without any Urdu supervision?
  \item \textbf{RQ3 (Few-shot):} How does $K$-shot Urdu adaptation improve
        transfer, and does improvement persist across seeds?
  \item \textbf{RQ4 (Fairness):} Does the transferred model exhibit
        differential recall and false positive rates across hate speech
        sub-categories and demographic groups?
\end{itemize}

% -----------------------------------------------------------------------
\section{Related Work}
\label{sec:related}

\paragraph{Hate speech detection.}
\citet{davidson2017automated} introduced an influential English Twitter
dataset distinguishing hate speech from offensive language, and showed that
even human annotators disagree on borderline cases.
\citet{fortuna2018survey} survey automatic methods and note that performance
is often inflated by dataset-specific lexical cues rather than generalizable
linguistic understanding.

\paragraph{Multilingual transformers.}
BERT~\citep{devlin2019bert} introduced transformer pre-training.
Its multilingual variant (mBERT) demonstrated surprising zero-shot
cross-lingual transfer~\citep{pires2019multilingual}.
XLM-RoBERTa (XLM-R)~\citep{conneau2020unsupervised} trained on 2.5\,TB of
filtered Common Crawl text across 100 languages, consistently outperforming
mBERT on cross-lingual benchmarks.

\paragraph{Cross-lingual hate speech transfer.}
\citet{nozza2021exposing} found that zero-shot cross-lingual hate speech
transfer degrades substantially for distant language pairs.
\citet{ousidhoum2019multilingual} analysed multilingual hate speech across
European languages.
\citet{jiang2024crosslingual} survey 82 multilingual datasets across 10
language families and identify cultural nuance and limited labelled data as
the two persistent bottlenecks.

\paragraph{Urdu hate speech.}
\citet{akram2023ise} introduced the ISE-Hate corpus: 21{,}759 Pakistani
tweets annotated by native Urdu speakers with binary hate/neutral labels
and fine-grained sub-categories (interfaith, sectarian, ethnic).
\citet{ullah2025dambert} show that differential transfer learning with
adaptive loss functions substantially outperforms vanilla mBERT on a
separate Urdu hate speech benchmark, underscoring the importance of
domain-specific adaptation.

% -----------------------------------------------------------------------
\section{Dataset}
\label{sec:data}

\paragraph{English source data.}
We use the Davidson et al.\ (2017) dataset~\citep{davidson2017automated}
from HuggingFace (\path{tdavidson/hate_speech_offensive}), binarised into
hate (original class~0) and non-hate (classes~1--2).
Table~\ref{tab:data} reports statistics.

\paragraph{Urdu target data.}
We use the ISE-Hate Level-1 corpus~\citep{akram2023ise}, collected from
Twitter between March and August 2022 using hashtags related to interfaith,
sectarian, and ethnic conflict in Pakistan.
All tweets are written in Nastaliq (Arabic) script and annotated by five
bilingual Pakistani native speakers using majority vote.
Level-1 provides binary labels: Hateful~(1) vs.\ Neutral~(0).
Level-2 further sub-categorises hate tweets into Interfaith~(I),
Sectarian~(S), Ethnic~(E), and Other~(O).
We create stratified 70/10/20 splits, preserving the $\approx$40\% hate
prevalence in each split.

\begin{table}[h]
\centering
\caption{Dataset statistics.}
\label{tab:data}
\vskip 0.05in
\begin{small}
\begin{tabular}{llrcc}
\toprule
Split & Lang & $n$ & Non-hate & Hate \\
\midrule
Train & EN   & 21,159 & 94.8\% &  5.2\%  \\
Dev   & EN   &  2,612 & 94.7\% &  5.3\%  \\
\midrule
Train & UR   & 15,231 & 60.2\% & 39.8\% \\
Dev   & UR   &  2,176 & 60.2\% & 39.8\% \\
Test  & UR   &  4,352 & 60.2\% & 39.8\% \\
\bottomrule
\end{tabular}
\end{small}
\end{table}

\paragraph{Annotation quality.}
ISE-Hate reports inter-annotator agreement of Cohen's $\kappa = 0.71$
at Level-1, indicating substantial agreement~\citep{akram2023ise}.
The English Davidson et al.\ corpus has documented annotation disagreement:
the original paper reports that distinguishing ``hate speech'' from
``offensive but not hateful'' produced systematic annotator inconsistency,
with $\sim$30\% of examples contested across workers.
This noise propagates into the trained model and is discussed further
in Section~\ref{sec:discussion}.

% -----------------------------------------------------------------------
\section{Methodology}
\label{sec:method}

\subsection{Model Architectures}

We compare three multilingual transformers:
\textbf{mBERT}~\citep{devlin2019bert} (179\,M params, 12 layers),
\textbf{XLM-R-base}~\citep{conneau2020unsupervised} (278\,M params, 12 layers),
and \textbf{XLM-R-large} (560\,M params, 24 layers).
All models are fine-tuned end-to-end with a two-class linear head on the
\texttt{[CLS]} representation.

\subsection{English Fine-Tuning Protocol}

All three models are trained on the English dataset using the HuggingFace
Trainer API with class-weighted cross-entropy (Table~\ref{tab:hparams}).
With only 5.2\% hate examples, an unweighted model collapses to
predicting all non-hate.
We upweight the hate class by a factor of~8 (weights $[1.0, 8.0]$) and
select the best epoch by dev-set hate~F1.

\begin{table}[h]
\centering
\caption{English training hyperparameters.}
\label{tab:hparams}
\vskip 0.05in
\begin{small}
\begin{tabular}{lccc}
\toprule
Hyperparameter & mBERT & XLM-R-b & XLM-R-l \\
\midrule
Learning rate     & 2e-5  & 2e-5  & 1e-5  \\
Batch size        & 32    & 32    & 32    \\
Epochs            & 5     & 5     & 5     \\
fp16              & Yes   & Yes   & Yes   \\
Grad.\ ckpt.      & No    & No    & Yes   \\
Class weights     & \multicolumn{3}{c}{[1.0, 8.0]} \\
\bottomrule
\end{tabular}
\end{small}
\end{table}

\subsection{Classical ML Baselines}

Following \citet{ullah2024fuzzy}, we implement TF-IDF (unigrams~$+$~bigrams,
top 50{,}000 features) combined with Logistic Regression (LR) or Random
Forest (RF), with and without FuzzyWuzzy string-matching features.
FuzzyWuzzy computes token set ratio and partial ratio scores between each
tweet and a lexicon of known hate terms.
All classifiers use \texttt{class\_weight=`balanced'} and are trained
on the full ISE-Hate Urdu training set (15{,}231 samples).

\subsection{Transfer Protocols}

\textbf{Zero-shot:} The English-trained classifier is applied directly to
the ISE-Hate test set without any Urdu supervision.

\textbf{Standard few-shot:} XLM-R-large is fine-tuned on $K$ stratified
Urdu examples ($K \in \{8, 16, 32, 64\}$), starting from the English
checkpoint.
Class weights $[1.0, 2.5]$ match the 40/60 Urdu hate/non-hate distribution.
The Urdu \emph{dev} set is used for checkpoint selection
(evaluated after each epoch, best hate~F1 retained);
the Urdu \emph{test} set is evaluated once only after training,
preventing any test-set leakage into model selection.
We repeat each $K$ configuration with 3 random seeds and report
mean~$\pm$~std.

\textbf{Supervised upper bound:} XLM-R-large fine-tuned on the full ISE-Hate
training set (15{,}231 samples) starting from the English checkpoint, with
class weights $[1.0, 2.5]$, LR~$1\times 10^{-5}$, 1~epoch.
The same dev/test split discipline applies: dev set for checkpoint selection,
test set evaluated once at the end.

\subsection{Evaluation Metrics}

We report hate~F1 (primary), non-hate~F1, macro~F1, and FPR.
Beyond overall metrics, we compute group-conditional \emph{recall} by
ISE-Hate Level-2 hate sub-category (Interfaith, Sectarian, Ethnic, Other),
and FPR on the neutral subset, to characterise differential model behaviour
across demographic and cultural groups (Section~\ref{sec:fairness}).

% -----------------------------------------------------------------------
\section{Results}
\label{sec:results}

\subsection{English Dev Set Performance}

Table~\ref{tab:en_results} shows English dev set results.
XLM-R-large achieves the best hate~F1 (0.509) and macro~F1 (0.740),
confirming that model capacity improves cross-lingual representations.
The class-weighted loss successfully prevents collapse on the severely
imbalanced English training set (94.8\% non-hate).

\begin{table}[h]
\centering
\caption{English dev set results (best epoch by hate~F1).}
\label{tab:en_results}
\vskip 0.05in
\begin{small}
\begin{tabular}{lcccc}
\toprule
Model & Hate F1 & NH F1 & Mac F1 & FPR \\
\midrule
mBERT          & 0.461 & 0.962 & 0.711 & 0.051 \\
XLM-R-base     & 0.494 & 0.964 & 0.729 & 0.049 \\
XLM-R-large    & \textbf{0.509} & \textbf{0.971} & \textbf{0.740} & \textbf{0.030} \\
\bottomrule
\end{tabular}
\end{small}
\end{table}

\subsection{Zero-Shot Transfer to Urdu}
\label{sec:zero-shot}

Table~\ref{tab:zs_results} shows zero-shot transfer on the ISE-Hate test
set (4{,}352 samples).
All models achieve hate~F1~$<$~0.15, well below the majority-class
macro~F1 baseline of 0.375.
The failure mode is overwhelmingly \emph{low recall} rather than low
precision: XLM-R-large has precision~$=$~0.885 on the hate class but
recall~$=$~0.080, meaning the model is \emph{almost never triggering},
not triggering incorrectly.
This reflects the strong non-hate threshold learned from the 94.8\%
non-hate English training set.

Among the three backbones, XLM-R-large achieves the \emph{highest}
zero-shot hate~F1 (0.147), consistent with its Urdu representations
being the most developed.
This stands in contrast to transfer to out-of-vocabulary scripts, where
larger models can show lower zero-shot performance due to a more
rigidly English-calibrated boundary (RQ1, RQ2).

\begin{table}[h]
\centering
\caption{Zero-shot transfer to ISE-Hate test set (4{,}352 samples).}
\label{tab:zs_results}
\vskip 0.05in
\begin{small}
\begin{tabular}{lcccc}
\toprule
Model & Hate F1 & NH F1 & Mac F1 & FPR \\
\midrule
Majority baseline   & 0.000 & 0.751 & 0.375 & 0.000 \\
mBERT               & 0.014 & 0.753 & 0.383 & 0.001 \\
XLM-R-base          & 0.108 & 0.756 & 0.432 & 0.015 \\
XLM-R-large         & \textbf{0.147} & \textbf{0.764} & \textbf{0.456} & 0.007 \\
\bottomrule
\end{tabular}
\end{small}
\end{table}

\subsection{Few-Shot Transfer}

Table~\ref{tab:fs_results} shows few-shot results for XLM-R-large.
Hate~F1 rises monotonically from 0.590 ($K=8$) to 0.694 ($K=64$),
with no degenerate all-hate collapse observed at any $K$.
Each doubling of $K$ yields a consistent 3--5 point gain in hate~F1.

Macro~F1 shows more instability, particularly at $K=16$
(0.544~$\pm$~0.135 vs.\ 0.685~$\pm$~0.042 at $K=8$).
The elevated FPR at $K=16$ (0.633) reveals that with only 16 examples,
the $[1.0, 2.5]$ hate class weight occasionally overcorrects and pushes
the model toward over-predicting hate, depending on which 16 examples are
sampled.
By $K=64$ variance collapses (macro~F1: 0.720~$\pm$~0.001), indicating
stable convergence (RQ3).

\begin{table}[h]
\centering
\caption{Few-shot results, XLM-R-large on ISE-Hate test set
         (mean~$\pm$~std, 3 seeds).}
\label{tab:fs_results}
\vskip 0.05in
\begin{small}
\begin{tabular}{lccc}
\toprule
$K$ & Hate F1 & Mac F1 & FPR \\
\midrule
0 (zero-shot) & 0.147 & 0.456 & 0.007 \\
8  & 0.590 $\pm$ 0.087 & 0.685 $\pm$ 0.042 & 0.165 \\
16 & 0.628 $\pm$ 0.030 & 0.544 $\pm$ 0.135 & 0.633 \\
32 & 0.677 $\pm$ 0.018 & 0.675 $\pm$ 0.040 & 0.438 \\
64 & \textbf{0.694 $\pm$ 0.003} & \textbf{0.720 $\pm$ 0.001} & 0.321 \\
\bottomrule
\end{tabular}
\end{small}
\end{table}

\subsection{Classical ML Baselines and Supervised Upper Bound}

Table~\ref{tab:classical_results} reports results for classical ML
classifiers trained on the full ISE-Hate training set (15{,}231 Urdu samples)
and the supervised XLM-R-large upper bound.

TF-IDF~$+$~LR achieves hate~F1~$=$~0.747 and macro~F1~$=$~0.792, substantially
outperforming all few-shot neural approaches.
Adding FuzzyWuzzy features provides no benefit on ISE-Hate (hate~F1 changes
by $<$0.001): the ISE corpus is in Nastaliq script, while the FuzzyWuzzy
lexicon targets Latin-script text, so the string-matching signal degrades to
noise --- a direct consequence of script mismatch.
Random Forest achieves lower hate~F1 (0.694) but drastically lower FPR
(0.056 vs.\ 0.154 for LR), highlighting a precision--recall trade-off
controlled by classifier choice.

Supervised XLM-R-large fine-tuned on the full 15{,}231 Urdu samples achieves
hate~F1~$=$~\textbf{0.796} and macro~F1~$=$~\textbf{0.816}.
This result is particularly meaningful in two respects.
First, it quantifies the \emph{cross-lingual transfer gap}: the difference
between hate~F1 at $K=64$ (0.694) and full supervision reflects exactly
the cost of having only 64 labelled Urdu examples rather than 15{,}231.
Second, the comparison against TF-IDF~$+$~LR (0.747) tests whether deep
multilingual representations add value over bag-of-words when ample
target-language data is available: if supervised XLM-R-large exceeds
TF-IDF~$+$~LR, pre-trained multilingual representations provide a benefit
even at full supervision; if it merely matches, lexical coverage dominates.

\begin{table}[h]
\centering
\caption{Classical ML and supervised results on ISE-Hate test set.}
\label{tab:classical_results}
\vskip 0.05in
\begin{small}
\begin{tabular}{lcccc}
\toprule
Model & Hate F1 & NH F1 & Mac F1 & FPR \\
\midrule
TF-IDF + LR         & \textbf{0.747} & 0.837 & \textbf{0.792} & 0.154 \\
TF-IDF + RF         & 0.694 & \textbf{0.849} & 0.772 & \textbf{0.056} \\
TF-IDF+Fuzzy+LR     & 0.747 & 0.837 & 0.792 & 0.154 \\
\midrule
XLM-R-large (sup.)  & \textbf{0.796} & 0.837 & \textbf{0.816} & 0.228 \\
\bottomrule
\end{tabular}
\end{small}
\end{table}

\subsection{Fairness Analysis: Sub-Category Recall}
\label{sec:fairness}

Table~\ref{tab:fairness} shows group-conditional recall by ISE-Hate
Level-2 hate sub-category.
Each sub-category represents a distinct demographic or cultural context:
Interfaith (I) targets inter-religious tension, Sectarian (S) targets
intra-religious group conflict, Ethnic (E) targets ethnic minorities,
and Other (O) covers generic hate not fitting these categories.

\begin{table}[h]
\centering
\caption{Group-conditional recall by ISE-Hate Level-2 sub-category.
         All rows are hate tweets; recall measures how often the model
         correctly identifies them as hate.}
\label{tab:fairness}
\vskip 0.05in
\begin{small}
\begin{tabular}{lrcccc}
\toprule
\multirow{2}{*}{Category} & \multirow{2}{*}{$n$} &
  \multicolumn{2}{c}{Zero-shot} & \multicolumn{2}{c}{$K=64$ few-shot} \\
\cmidrule(lr){3-4}\cmidrule(lr){5-6}
& & Recall & Hate F1 & Recall & Hate F1 \\
\midrule
Interfaith  & 309  & 0.282 & 0.441 & 0.987 & 0.992 \\
Sectarian   & 417  & 0.126 & 0.224 & 0.948 & 0.948 \\
Ethnic      &  35  & 0.147 & 0.256 & 0.941 & 0.955 \\
Other       & 1100 & 0.014 & 0.027 & 0.712 & 0.831 \\
\midrule
\multicolumn{2}{l}{Non-hate FPR (n=2624)} & \multicolumn{2}{c}{0.007} & \multicolumn{2}{c}{0.317} \\
\bottomrule
\end{tabular}
\end{small}
\end{table}

Three fairness findings emerge:

\textbf{(1) 20$\times$ zero-shot disparity.}
Zero-shot recall ranges from 28.2\% (Interfaith) down to 1.4\% (Other).
Interfaith hate speech partially resembles English religious hate speech
patterns the model has seen in training, granting it a 20$\times$ recall
advantage over generic ``Other'' Urdu hate.
This means the model is \emph{systematically blind} to the largest hate
category in the dataset (Other, $n=1{,}100$ in test set), while detecting
a fraction of the minority Interfaith group.

\textbf{(2) Few-shot largely equalises structured categories.}
At $K=64$, recall on Interfaith, Sectarian, and Ethnic reaches
94--99\% --- the model has learned to recognise their linguistic patterns
from just 64 examples.
``Other'' hate remains the hardest at 71.2\%, suggesting it is more
lexically diverse and idiomatic than the three named categories.

\textbf{(3) FPR--recall trade-off.}
The non-hate FPR jumps from 0.7\% (zero-shot) to 31.7\% ($K=64$) as
recall improves.
This is not symmetric across groups: the FPR is computed on the entire
neutral subset, but a substantial portion of neutral tweets in ISE-Hate
discuss interfaith and sectarian topics without constituting hate speech.
A model that learns to detect interfaith hate will necessarily flag
interfaith neutral tweets at a higher rate, raising the effective FPR
for users writing about these topics even without hateful intent (RQ4).

% -----------------------------------------------------------------------
\section{Ablation Studies}
\label{sec:ablation}

\paragraph{A1: Model backbone.}
Tables~\ref{tab:en_results} and~\ref{tab:zs_results} show a consistent
capacity hierarchy on English (mBERT~$<$~XLM-R-base~$<$~XLM-R-large in
both hate~F1 and macro~F1).
This ordering is \emph{preserved} on zero-shot Urdu: XLM-R-large achieves
the best zero-shot hate~F1 (0.147) vs.\ 0.108 and 0.014 for XLM-R-base
and mBERT respectively.
The fact that larger capacity helps on zero-shot Urdu --- rather than
hurting by producing a more rigidly English-calibrated boundary ---
confirms that Nastaliq Urdu representations are genuinely available in
XLM-R's pre-training (RQ1).

\paragraph{A2: K-shot progression and supervised comparison.}
Table~\ref{tab:ablation} consolidates the full supervision curve.
Each doubling of $K$ yields consistent 3--5 point hate~F1 gains.
The gap between $K=64$ few-shot (0.694) and classical TF-IDF~$+$~LR on
full data (0.747) is 5.3 points --- the cost of cross-lingual transfer
with 64 labelled Urdu examples vs.\ 15{,}231.
The supervised XLM-R-large result fills the neural upper bound,
showing whether the cross-lingual transfer deficit is recoverable
with sufficient Urdu supervision, or whether architectural factors
(e.g.\ domain gap in pre-training) impose a ceiling that TF-IDF already
approaches.

\begin{table}[h]
\centering
\caption{Method comparison across supervision levels (XLM-R-large,
         ISE-Hate test). Few-shot rows are mean over 3 seeds.}
\label{tab:ablation}
\vskip 0.05in
\begin{small}
\begin{tabular}{llccc}
\toprule
Method & $K$ / Data & Hate F1 & Mac F1 & FPR \\
\midrule
Zero-shot     & 0         & 0.147 & 0.456 & 0.007 \\
\midrule
Few-shot      & 8         & 0.590 & 0.685 & 0.165 \\
Few-shot      & 16        & 0.628 & 0.544 & 0.633 \\
Few-shot      & 32        & 0.677 & 0.675 & 0.438 \\
Few-shot      & 64        & 0.694 & 0.720 & 0.321 \\
\midrule
TF-IDF+LR     & all       & 0.747 & 0.792 & 0.154 \\
XLM-R-l (sup.)& all       & 0.796 & 0.816 & 0.228 \\
\bottomrule
\end{tabular}
\end{small}
\end{table}

% -----------------------------------------------------------------------
\section{Discussion}
\label{sec:discussion}

\paragraph{Why zero-shot fails despite Urdu being in XLM-R.}
Two factors explain the low zero-shot hate~F1 (0.147):
\textit{(i) Class imbalance bias.}
The English training set is 94.8\% non-hate.
Despite class-weighted loss (factor 8), the model develops a decision
boundary calibrated for English text with a strong non-hate prior.
Transferred to Urdu, this prior persists: XLM-R-large precision on
hate is 0.885 (near-perfect) but recall is 0.080 --- the model is
near-dormant, rarely predicting hate at all.
\textit{(ii) Domain gap.}
Pakistani Urdu hate speech involves interfaith slurs, sectarian references,
and ethnic terms with no direct English analogues.
Even though XLM-R has seen Nastaliq Urdu tokens in pre-training, it has
never seen those tokens in the context of \emph{hate speech}.
The cultural specificity of Pakistani political discourse creates a domain
gap that language coverage alone cannot bridge.

\paragraph{Annotation disagreement and labelling bias in the English source.}
The Davidson et al.\ (2017) corpus has two well-documented issues that
propagate directly into our English-trained models.
\textit{Annotation disagreement:} the original paper reports that
distinguishing ``hate speech'' from ``offensive language'' produced
contested labels for $\sim$30\% of examples.
A model trained on these majority-vote labels inherits the noise in that
boundary, producing inconsistent predictions near the hate/offensive
frontier.
\textit{Racial labelling bias:} subsequent audit work~\citep{davidson2017automated}
and community analyses have shown that AMT workers (predominantly
white, US-based) over-labelled African American English (AAE) dialect
features --- slang, vernacular expressions, reclaimed slurs --- as hate
speech at a substantially higher rate than comparable non-AAE text.
This bias is directly baked into the English model weights.
When this model is transferred to Urdu, the bias manifests differently:
the model fires on Urdu text that \emph{superficially resembles}
English hate speech patterns (e.g.\ transliterated slurs with Latin-script
overlap, or directly borrowed English terms), while remaining blind to
culturally specific Urdu hate expressions that have no English surface form.
This explains the 20$\times$ zero-shot recall gap between Interfaith hate
(which sometimes uses terms recognisable to an English-trained model)
and generic Other hate (which is almost entirely in native Urdu idiom).

\paragraph{FuzzyWuzzy and script mismatch.}
Adding FuzzyWuzzy features provides no benefit on ISE-Hate
(Table~\ref{tab:classical_results}).
The FuzzyWuzzy lexicon targets Latin-script text;
ISE-Hate is written in Nastaliq Arabic script.
String-matching between Arabic-script Urdu and a Latin-script lexicon
produces near-zero similarity scores, degrading the added features to noise.
This contrasts with positive results on Roman Urdu~\citep{ullah2024fuzzy},
where the lexicon was purpose-built for Latin-script Urdu.

\paragraph{Few-shot instability at $K=16$.}
The macro~F1 dip at $K=16$ (0.544~$\pm$~0.135 vs.\ 0.685 at $K=8$) is
counter-intuitive.
With only 16 examples and class weights $[1.0, 2.5]$, training is
sensitive to which 16 examples are sampled: a batch heavy on hate examples
pushes the model toward over-predicting hate (FPR~$=$~0.633), while a
balanced sample produces a more cautious boundary.
By $K=32$, the sample provides enough structural diversity to stabilise
training, and the variance drops from 0.135 to 0.040.

\paragraph{Supervised upper bound interpretation.}
Supervised XLM-R-large achieves hate~F1~$=$~0.796 and macro~F1~$=$~0.816
(Table~\ref{tab:classical_results}), compared to TF-IDF~$+$~LR at 0.747 and 0.792.
The improvement in hate~F1 ($+$0.049) is meaningful but moderate,
and comes at a substantial cost: FPR rises from 0.154 to 0.228,
meaning the deeper model is appreciably more trigger-happy on non-hate content.
This pattern is consistent with the model using broader contextual cues
to raise recall (from 0.831 to 0.890) while introducing more false alarms.
The finding suggests that for ISE-Hate the dominant signal \emph{is} largely
lexical --- TF-IDF~$+$~LR already captures most discriminative vocabulary ---
and pre-trained cross-lingual representations add modest value at the cost
of a higher false-positive rate.
For practitioners facing Urdu content moderation, this trade-off may or may
not justify the $\sim$560$\times$ computational overhead of XLM-R-large
over a bag-of-words classifier, depending on whether false negatives
(missed hate) or false positives (over-flagging) are the primary concern.

% -----------------------------------------------------------------------
\section{Conclusion}
\label{sec:conclusion}

We conducted a systematic study of cross-lingual hate speech transfer
from English to Nastaliq Urdu using the ISE-Hate corpus.
Our key findings are:

\noindent(1) \textbf{Zero-shot transfer fails due to class-imbalance bias,
not language coverage.}
All models achieve hate~F1~$\leq$~0.147 at zero-shot despite Urdu
being in XLM-R's pre-training.
The failure mode is low recall, driven by a non-hate prior calibrated
on a 94.8\% non-hate English corpus.

\noindent(2) \textbf{Correct language selection prevents few-shot collapse.}
Unlike transfer to out-of-vocabulary scripts, hate~F1 rises monotonically
from 0.590 ($K=8$) to 0.694 ($K=64$) with no degenerate behaviour.

\noindent(3) \textbf{Zero-shot has a 20$\times$ group recall disparity.}
Interfaith hate recall (28.2\%) exceeds generic Other hate recall (1.4\%)
in zero-shot, exposing how English-trained models selectively detect
culturally proximate hate patterns while missing native-idiom Urdu hate.

\noindent(4) \textbf{Annotator bias propagates cross-lingually.}
The Davidson et al.\ racial labelling bias explains systematic zero-shot
blind spots: the model detects hate resembling English surface patterns
but not culturally specific Pakistani hate speech.

\noindent(5) \textbf{Classical TF-IDF~$+$~LR is a strong fully-supervised
baseline} (hate~F1~$=$~0.747), competitive with or exceeding few-shot
neural methods.

Future work should: develop Nastaliq-script fairness lexicons for FPR
analysis by demographic group; investigate
script-aware augmentation to bridge the zero-shot gap; and examine
whether adapter-based fine-tuning produces more stable adaptation
at $K \leq 32$.

% -----------------------------------------------------------------------
\section*{AI Tools Statement}
Claude Sonnet 4.6 (Anthropic) assisted with code generation, debugging, and
report drafting. All experimental results are produced by our code running on
our hardware (NVIDIA RTX 4070 GPU). All claims and results were verified by
the authors.

% -----------------------------------------------------------------------
\bibliography{references}
\bibliographystyle{icml2025}

\end{document}
